\documentclass[11pt]{article}
\usepackage[margin=1in]{geometry}
\usepackage{amsmath,amssymb}
\usepackage{hyperref}

\title{G8 Lane A Strong Standard-Mode Proof Attempt and Pivot}
\author{Private red-team proof-design memo}
\date{May 21, 2026}

\newcommand{\Spec}{\operatorname{Spec}}
\newcommand{\im}{\operatorname{Im}}
\newcommand{\re}{\operatorname{Re}}

\begin{document}
\maketitle

\section*{Red-Team Verdict}

The strong exact standard-mode theorem is now isolated with very high
precision.  In mathematical form it says:
\[
  \forall z\in \mathcal Z_{\xi}^{\neq 0},\quad
  \exists n\in\mathbb N,\quad
  V(z)=\lambda_n,
\]
where
\[
  V(z)=\iota_{\tau,500}^2
  \left(s_z(1-s_z)-\frac14\right)
\]
is the Lean-selected canonical scaled actual-\(\xi\) readout, and
\[
  \lambda_n =
  \texttt{Tau.BookIII.Doors.lemniscate\_eigenvalue}\ n
  =
  n^2
\]
in the current standard lemniscate eigenvalue model.

This theorem is stronger than Lane A strictly needs.  It implies spectral
reality, but it also asserts exact quantization into the present \(n^2\) ladder:
for a critical-line point \(s_z=\frac12+i\gamma_z\), it amounts morally to
\[
  \iota_{\tau,500}^2\gamma_z^2=n^2.
\]
The current Lean and manuscript sources do not yet prove this global
quantization theorem for every actual nonzero-height \(\xi\)-carrier.

Therefore the strong theorem should be treated as a conditional reduction unless
a new independent spectral-coordinate theorem is supplied.  The weaker abstract
spectral-membership route is the better next target for Lane A, because it asks
only for genuine self-adjoint spectral membership and real-valuedness, not exact
membership in the current \(n^2\) enumeration.

\section*{Lean Targets}

The strong theorem is represented by:
\[
\begin{aligned}
&\texttt{modeOf : G8ActualXiCanonicalModeExtractionTarget},\\
&\texttt{alignment :
G8ActualXiCanonicalModeEigenvalueAlignmentTarget modeOf}.
\end{aligned}
\]
The alignment field expands to:
\[
  \forall z,\quad
  V(z)
  =
  \left(
  \texttt{Tau.BookIII.Doors.lemniscate\_eigenvalue}\ (\texttt{modeOf}\ z)
  :\mathbb C\right).
\]
Once supplied, this packages as
\[
  \texttt{G8ActualXiCanonicalStandardEigenvalueMembershipSource}.
\]

The weaker source-facing target is:
\[
  \texttt{G8ActualXiIotaTauCanonicalSelfAdjointRealityTarget}.
\]
Even more explicitly, the Lane-A-sufficient payload is:
\[
  \forall z\in\mathcal Z_\xi^{\neq 0},\quad
  \im(V(z))=0.
\]

\section*{Strong Proof Attempt}

\paragraph{Theorem A: Exact standard-mode source.}
For every actual nonzero-height \(\xi\)-carrier \(z\), there exists a natural
number \(n(z)\) such that
\[
  V(z)=\lambda_{n(z)}.
\]

\paragraph{Attempted proof.}
\begin{enumerate}
\item Fix \(z\in\mathcal Z_\xi^{\neq 0}\), with orthodox point \(s_z\).
\item By the canonical scaled image law,
  \[
    V(z)=\iota_{\tau,500}^2
    \left(s_z(1-s_z)-\frac14\right).
  \]
\item Prove that \(V(z)\) is a genuine spectral value of the ready Book III
  lemniscate operator:
  \[
    V(z)\in \Spec(H_L).
  \]
\item Prove a standard spectrum classification theorem:
  every member of the chosen standard lemniscate spectrum is exactly
  \[
    \lambda_n=n^2
  \]
  for some \(n\in\mathbb N\), with the same normalization as the Lean
  definition \(\texttt{lemniscate\_eigenvalue}\).
\item Choose \(n(z)\) from this classification.
\item Define \(\texttt{modeOf}\ z:=n(z)\), obtaining the exact alignment field.
\end{enumerate}

\paragraph{First blocked lemma.}
The first missing lemma is:
\[
  \forall z\in\mathcal Z_\xi^{\neq 0},\quad
  V(z)\in \Spec_{\mathrm{std}}(H_L),
\]
where
\[
  \Spec_{\mathrm{std}}(H_L)=\{n^2:n\in\mathbb N\}
\]
is the current standard Lean spectrum.

This is not supplied by self-adjointness.  Self-adjointness proves that spectral
members are real; it does not prove that a given actual-\(\xi\) value is a
member, and it does not identify that value with an integer square.

\section*{Why The Strong Target Is Too Strong Right Now}

\paragraph{Integer-square quantization.}
The current standard eigenvalue model is not merely a generic real spectrum.
It is the discrete ladder \(\lambda_n=n^2\).  Thus the strong theorem requires:
\[
  \iota_{\tau,500}^2
  \left(s_z(1-s_z)-\frac14\right)=n^2.
\]
After the critical-axis simplification this becomes the exact height
quantization condition:
\[
  \iota_{\tau,500}\,|\gamma_z|=n.
\]
No current source theorem proves such a global exact quantization law for the
actual \(\xi\)-zero stream.

\paragraph{Manuscript support.}
Book III Chapter 23 supports the self-adjoint metric-graph operator and its
real, discrete spectrum.  Chapter 24 defines
\[
  \Lambda(s)=\iota_\tau^2\left(s(1-s)-\frac14\right)
\]
and routes zero-to-spectrum membership through the determinant/spectral
correspondence bridge.  Chapter 25 uses only spectral reality to force the
critical line.  It does not require exact equality to \(n^2\).

\paragraph{Finite diagnostics.}
Finite projector checks, finite spectral-resolution checks, finite
correspondence checks, and finite self-adjoint checks are valuable diagnostics.
They do not construct a global mode function \(z\mapsto n(z)\), and they do not
prove exact equality \(V(z)=n^2\).

\paragraph{Normalization risk.}
The Lean value \(\iota_{\tau,500}^2\) is the certified 500-digit tower scale.
It is positive and nonzero, which is exactly what the spectral-reality route
needs.  But the strong standard-mode theorem needs more: it needs exact
compatibility between this selected scale, the operator normalization, and the
integer-square enumeration.

\section*{What The Strong Theorem Would Give}

If Theorem A were supplied, the rest of Lane A is already a proof-checked
adapter chain:
\begin{enumerate}
\item Exact standard membership gives a standard eigenvalue index source.
\item Standard eigenvalues are real-valued after complex readout.
\item Therefore \(\im(V(z))=0\).
\item Since \(\iota_{\tau,500}^2\neq0\), the unscaled centered quadratic has
  zero imaginary coordinate.
\item The algebraic identity
  \[
    \im\left(s_z(1-s_z)-\frac14\right)
    =
    \gamma_z(1-2\sigma_z)
  \]
  is already formalized.
\item Since \(z\) is nonzero-height, \(\gamma_z\neq0\), hence
  \(\sigma_z=\frac12\).
\item Combined with the closed zero-height eta branch, this gives actual
  sigma-fixedness.
\end{enumerate}

\section*{Recommended Pivot}

\paragraph{Theorem B: Abstract self-adjoint spectral membership.}
Let \(S\) be a genuine self-adjoint lemniscate spectral source with membership
predicate \(\operatorname{Mem}_S(\lambda)\).  Assume:
\[
  \operatorname{Mem}_S(\lambda)\Longrightarrow \im(\lambda)=0
\]
and
\[
  \forall z\in\mathcal Z_\xi^{\neq0},\quad \operatorname{Mem}_S(V(z)).
\]
Then:
\[
  \forall z\in\mathcal Z_\xi^{\neq0},\quad
  \im\left(s_z(1-s_z)-\frac14\right)=0.
\]
Thus the nonzero-height Lane A spectral-parameter reality theorem follows.

\paragraph{Proof of Theorem B.}
\begin{enumerate}
\item Fix \(z\in\mathcal Z_\xi^{\neq0}\).
\item By canonical membership, \(\operatorname{Mem}_S(V(z))\).
\item By self-adjoint spectral reality, \(\im(V(z))=0\).
\item By the canonical scaled image law,
  \[
    V(z)=\iota_{\tau,500}^2
    \left(s_z(1-s_z)-\frac14\right).
  \]
\item Since \(\iota_{\tau,500}^2\) is a nonzero real scalar, the centered
  quadratic has zero imaginary coordinate.
\item The centered quadratic identity and nonzero-height condition force
  \(\re(s_z)=1/2\).
\end{enumerate}

\section*{Why The Pivot Is Enough For Lane A}

Lane A does not need exact mode extraction.  It needs real-valuedness of the
canonical scaled spectral value.  Exact standard-mode membership is one way to
obtain that real-valuedness, but it is not the minimal way.  The minimal
operator theorem is:
\[
  \forall z,\quad \im(V(z))=0,
\]
preferably sourced as genuine membership in a self-adjoint lemniscate spectrum.

This aligns with Book III Chapter 25: the critical-line forcing argument uses
spectral reality, not integer-square mode enumeration.

\section*{Next Lean Translation Target}

The recommended next Lean module should add the weaker source surface before
attempting the standard-mode ladder again:
\[
  \texttt{G8ActualXiCanonicalAbstractSpectralMembershipSource}.
\]
It should carry:
\begin{enumerate}
\item a genuine spectral membership predicate for the ready lemniscate operator;
\item a self-adjoint real-valuedness selector for that predicate;
\item pointwise membership of \(V(z)\) for every actual nonzero-height carrier;
\item adapters to \(\texttt{G8ActualXiNonzeroHeightSpectralParameterReal}\).
\end{enumerate}

The strong standard-mode theorem can then be pursued later as a refinement:
\[
  \operatorname{Mem}_S(V(z))
  \quad+\quad
  \text{standard spectrum classification}
  \quad\Longrightarrow\quad
  \exists n,\ V(z)=n^2.
\]
This separates the Lane-A proof need from the additional exact enumeration and
scale-calibration burden.

\end{document}
